%% kapa: exam=1 round=34 year=2023 period=1 form=A subject=economics subject_ko=경제학원론 seq=11 number=51
%% source: Q-Net 감정평가사 기출문제 (cst00309) artlSeq=5211789 file=raw/1008/5211789/2023년도 제34회 감정평가사 1차시험 1교시.pdf
%% conversion: kordoc 4.24.51 -> tools/gen1.py (auto); review=proofread-sonnet; flags=none
\begin{question}{51}
복점시장에서 기업 1과 기업 2는 각각 a와 b의 전략을 갖고 있다. 성과 보수 행렬이 다음과 같을 때, 내쉬균형을 모두 고른 것은? (단, 보수 행렬 내 괄호 안 왼쪽은 기업 1의 보수, 오른쪽은 기업 2의 보수이다.)

{\small\centering\renewcommand{\arraystretch}{1.3}
\begin{tabular}{|c|c|c|c|}\hline
\multicolumn{2}{|c|}{\rule{0pt}{1.2em}} & \multicolumn{2}{c|}{기업 2} \\ \cline{3-4}
\multicolumn{2}{|c|}{} & 전략 a & 전략 b \\ \hline
\multirow{2}{*}{기업 1} & 전략 a & (16, 8) & (8, 6) \\ \cline{2-4}
 & 전략 b & (3, 7) & (10, 11) \\ \hline
\end{tabular}\par}
\begin{choices}
  \item (16, 8)
  \item (10, 11)
  \item (8, 6), (10, 11)
  \item (16, 8), (3, 7)
  \item (16, 8), (10, 11)
\end{choices}
\end{question}
\answer{5}
